\begin{afterword}
\summaryhead

The author separates two early steps from the ``Bayesian enlightenment'' proper. The first is an anecdote. Asked about a mathematician who says ``I have two children, and at least one of them is a boy,'' the author answered that the chance of two boys is 1/2, since a mother of a boy and a girl could as well have mentioned the girl. Someone replied that this was the Bayesian answer and that orthodox statistics gives 1/3. The second step was the heuristics-and-biases literature.

The enlightenment came from E.~T. Jaynes's \textsc{Probability Theory: The Logic of Science}, which presented probability as rules that give ``the answer,'' not ``an answer.'' The author then held Friendly AI designs to the standard of a passage: ``do only that which you must do, and which you cannot do in any other way.'' The old theories were discarded for the study of probability and decision theory. Later, reading Judea Pearl, the author found that precision saves time. The post ends by calling the scream of embarrassment at a past self's mistakes ``the sound that rationalists make when they level up.''

\responsehead

The post is memoir, and the parts that can be checked mostly hold. The arithmetic of the puzzle is right. The author's objection, that the answer 1/3 assumes the mother would always mention a boy, is the standard analysis of this puzzle since Martin Gardner granted that his version was ambiguous. The description of Jaynes matches Jaynes's own claims: rules that do ``uniquely and optimally'' what orthodox statistics does with ``intuitive ad hoc devices,'' and results checked by computing them in several ways.

The frame of the anecdote does not hold. The difference between 1/3 and 1/2 lies in how the statement came to be made, not in the school of the person computing. A frequentist who counts the families in which a parent would say this sentence also gets 1/2. The post takes its label from a person who tied the difference to the schools. The notes on ``Probability is in the Mind'' made the same point about puzzles of this kind.

The two comparisons with frequentist methods, of which ``no two of them agreed'' and whose calculations ``seemed like good ideas,'' repeat Jaynes's side of a dispute. Frequentist theory derives its methods from stated criteria and proves results about them, as in the Neyman--Pearson lemma. The fuller discussion, including in what sense the exact Bayesian calculation is optimal and how far Cox's theorem makes the rules unique, is in the notes on ``Beautiful Probability''.

One claim about other people is asserted without evidence: that they call precision impossible ``because human beings are lazy.'' Its support is two links, one of which draws its ``five minutes'' from a film scene.

In short: a memoir whose account of Jaynes matches Jaynes's own claims and whose puzzle is solved correctly, but whose puzzle does not separate Bayesians from frequentists, and whose remarks on frequentist methods repeat Jaynes's side of a dispute treated earlier in the book.
\end{afterword}
